
\subsection{Non-diagonal Three-Qubit Entangler}
\label{sec:results_nv_xzz}

\noindent
We now apply the interaction-resolved control framework of Sec.~\ref{sec:control_objective_phase_invariants} to synthesize a non-diagonal three-qubit entangling gate by performing the optimization in a diagonalizing frame. The target operation is
\begin{equation}
\label{eq:desired_u_nondiag}
U_{\mathrm{target}} = e^{\,i\frac{\pi}{4}X\!\otimes\!Z\!\otimes\!Z},
\end{equation}
which also generates a GHZ state up to local phases and thus produces genuine tripartite entanglement~\cite{coffman2000distributed,dur2000three}. This target gate is not naturally diagonal and therefore provides an explicit example of the scheme for a non-diagonal Pauli-string interaction. \\

% \medskip
% \noindent
% As before, we simulate a solid-state register consisting of an NV electronic spin ($A$) coupled to two proximal $^{13}$C nuclear spins ($B$, $C$), with the $^{14}$N nuclear spin fixed in a spectator $m_I$ manifold. The system dynamics are described by the effective
% Hamiltonian in Eq.~\eqref{eq:H_NV_RWA_main}, and the microwave control fields are parameterized according to Eq.~\eqref{eq:pulse_parametrization_example}. Within the logical
% $A\!\otimes\!B\!\otimes\!C$ subspace, we therefore consider the projected propagator $U^{A\otimes B\otimes C}_{(\mathrm{pulse})}(\mathbf p,T)$, corresponding to the restriction of the full propagator $U_{\mathrm{pulse}}(\mathbf p,T)$ defined in Eq.~\eqref{eq:U_pulse_definition} to the logical $A\!\otimes\!B\!\otimes\!C$ subspace. The optimization then seeks a diagonal operation in the Hadamard frame whose interaction content consists of a finite three-body phase and vanishing pairwise couplings.

\noindent
As before, we simulate a solid-state register consisting of an NV electronic spin ($A$) coupled to two proximal $^{13}$C nuclear spins ($B$, $C$), with the $^{14}$N nuclear spin fixed in a spectator $m_I$ manifold. The system dynamics are described by the effective Hamiltonian in Eq.~\eqref{eq:H_NV_RWA_main}, and the microwave control fields are parameterized according to Eq.~\eqref{eq:pulse_parametrization_example}. We therefore again consider the projected propagator $U^{A\otimes B\otimes C}_{(\mathrm{pulse})}(\mathbf p,T)$, corresponding to the restriction of the full propagator $U_{\mathrm{pulse}}(\mathbf p,T)$ defined in Eq.~\eqref{eq:U_pulse_definition} to the logical $A\!\otimes\!B\!\otimes\!C$ subspace.\\


\noindent
The interaction coordinates for this target gate are constructed in the \emph{Hadamard-sandwiched frame},
\begin{equation}
\label{eq:frame_transformation}
\begin{aligned}
U_{\mathrm{diag}} &= (H_A\!\otimes\!I_B\!\otimes\!I_C)\,
U^{A\otimes B\otimes C}_{(\mathrm{pulse})}(\mathbf p,T)\,
(H_A\!\otimes\!I_B\!\otimes\!I_C) \\
\end{aligned}
\end{equation}
where the target unitary is diagonal and the propagator can therefore be analyzed directly through its computational-basis phases and corresponding interaction coordinates, as given in Secs.~\ref{sec:interaction_resolved_representation}–\ref{sec:control_objective_phase_invariants}. After optimization, the Hadamard gates can be removed such that the same physical pulse implements the desired non-diagonal action.


\medskip
\noindent
In the Hadamard-sandwiched frame~\eqref{eq:frame_transformation}, the optimization proceeds as for the previous target gate, using the interaction coordinates introduced in Sec.~\ref{sec:interaction_resolved_representation}. The control objective specifies the target interaction phases
\begin{equation}
\begin{aligned}
&\phi^\star(\{a,b,c\}) = \tfrac{\pi}{4}, \\
&\phi^\star(\{a,b\})
=
\phi^\star(\{a,c\})
=
\phi^\star(\{b,c\})
= 0,
\end{aligned}
\label{eq:Uxzzinvariants}
\end{equation}
with the optimization minimizing the $\pi$-periodic cost function $\mathcal{J}_{\mathrm{3q}}$ defined in Eq.~\eqref{eq:phase-invariant-cost}, while the weights are chosen as $w(\{{a,b,c}\})=0.5$, $w(|S|=2)=0.2$, and $w(|S|=1)=0.0$. Explicit formulas for the interaction coordinates in terms of the computational-basis phase values $\{\phi_x\}_{x\in\{0,1\}^3}$ are given in Appendix~\ref{app:explicit_3q_invariants}. Using these interaction coordinates, we optimize a control pulse toward the same prescribed many-body interaction structure as for the diagonal target, expressed now in the Hadamard-sandwiched frame of Eq.~\eqref{eq:frame_transformation}.



\medskip
\noindent
After optimization, we return to the computational basis by discarding the Hadamard gates in Eq.~\eqref{eq:frame_transformation}, such that the effective action of the pulse is
\begin{equation}
U^{A\otimes B\otimes C}_{(\mathrm{pulse})}(\mathbf p,T)
\;\approx\;
e^{i\frac{\pi}{4}X\!\otimes\!Z\!\otimes\!Z}
\end{equation}
up to local transformations.

\medskip
\noindent
The optimized microwave envelope (Fig.~\ref{fig:nv_3q_pulse_xzz}) exhibits smooth amplitude and phase modulation over $T=1250$~ns, and consists of $N_c = 11$ frequency terms, as given in Eq.~\eqref{eq:pulse_parametrization}. The used parameters are provided in the public repository~\cite{baran2026_multiqubit_control} in the form of a parameter vector as in Eq.~\eqref{eq:pulse_parametrization_example}. It realizes the control objective in the diagonalizing (Hadamard) frame as a single shaped pulse. After the removal of the diagonalizing transformations, the same physical pulse implements the target non-diagonal entangler in the computational basis.

\begin{figure}[h!]
\centering
\includegraphics[width=0.47\textwidth]{Bilder/drives_concatenated_XZZ.png}
\caption{
Optimized time-modulated Rabi-frequency $\Omega(t)$ implementing the non-diagonal three-qubit entangler $e^{i(\pi/4)X\!\otimes\!Z\!\otimes\!Z}$ in the computational basis, synthesized via interaction-resolved optimization in the diagonalizing (Hadamard) frame.
}
\label{fig:nv_3q_pulse_xzz}
\end{figure}

\noindent
Next we validate the interaction dynamics by tracking the time evolution of the interaction coordinates $\phi_t(S)$ for all nonempty subsets $S\subseteq\{a,b,c\}$. This reveals that the control pulse builds up the desired three-body interaction phase $\phi_t(\{a,b,c\})$ in the diagonalizing frame while guiding the pairwise phase components $\phi_t(\{a,b\})$, $\phi_t(\{a,c\})$, and $\phi_t(\{b,c\})$ to be suppressed toward zero, up to the inherent $2\pi$-periodicity~\eqref{eq:pi_peridicity}. Figure~\ref{fig:nv_3q_invariants_vs_time_xzz} shows the resulting interaction-coordinate trajectories.

\begin{figure}[h!]
\centering
\includegraphics[width=0.49\textwidth]{Bilder/all_invariants_one_plot_XZZ.png}
\caption{
Time evolution of the interaction coordinates $\phi_t(S)$ extracted from the phase map in the diagonalizing (Hadamard-sandwiched) frame during the optimized non-diagonal gate pulse. As in the canonically diagonal case, in Fig.~\ref{fig:nv_3q_invariants_vs_time_ZZZ}, the three-body interaction coordinate $\phi_t(\{a,b,c\})$ (thick black line) converges to $\pi/4$, realizing the desired three-body interaction. The two-body interaction coordinates $\phi_t(\{a,b\})$, $\phi_t(\{a,c\})$, and $\phi_t(\{b,c\})$ directly go to zero in this case, corresponding to vanishing pairwise interaction content.
}
\label{fig:nv_3q_invariants_vs_time_xzz}
\end{figure}

\medskip
\noindent
Residual local phases can be assessed and corrected via the one-body interaction coordinates in the diagonalizing frame and yield the local correction operator
\begin{equation}
U_{\mathrm{loc}}^{(\mathrm{diag})}
=
\exp\!\Bigl(-i\sum_{j\in\{a,b,c\}}\phi(\{j\})\, Z_j\Bigr).
\end{equation}
Transforming back to the computational basis gives
\begin{equation}
\begin{aligned}
U_{\mathrm{loc}}^{(\mathrm{comp})}
&=
(H_A\!\otimes\!I_B\!\otimes\!I_C)\,
U_{\mathrm{loc}}^{(\mathrm{diag})}\,
(H_A\!\otimes\!I_B\!\otimes\!I_C) \\
&=
e^{i(\phi(\{a\}) X_A + \phi(\{b\}) Z_B + \phi(\{c\}) Z_C)} .
\end{aligned}
\end{equation}
Note that here the local $X$ rotation cannot be corrected virtually and an additional single-qubit gate is needed. The final operation is then
\begin{equation}
U_{\mathrm{final}}
=
U_{\mathrm{loc}}^{(\mathrm{comp})}\,U^{A\otimes B\otimes C}_{(\mathrm{pulse})}(\mathbf p,T).
\end{equation}

\medskip
\noindent
The fidelity with the target operation in Eq.~\eqref{eq:desired_u_nondiag} evaluates to
\begin{equation}
F_{\mathrm{XZZ}}
=
\frac{\big|\mathrm{Tr}\!\left[
\left(U_{\mathrm{target}}\right)^{\!\dagger}
U_{\mathrm{final}}
\right]\big|^{2}}{8^{2}}
= 0.9985,
\end{equation}
demonstrating that interaction-resolved optimization in a diagonalizing frame extends naturally to high-fidelity non-diagonal multiqubit gates.

\medskip
\noindent
This example demonstrates that the interaction-resolved control objectives apply also to non-diagonal multi-qubit gates. By performing the optimization in a suitable diagonalizing frame, the interaction structure of a Pauli-string generator can be accessed through the same interaction coordinates. 

